% Part: computability
% Chapter: recursive-functions
% Section: other-recursions

\documentclass[../../../include/open-logic-section]{subfiles}

\begin{document}

\olfileid{cmp}{rec}{ore}
\olsection{ఇతర పునరావృత్తులు}

జతీకరణ, క్రమీకరణలను ఉపయోగించి మరింత అసాధారణమైన (ఉపయోగకరమైన)
ఆదిమ పునరావృత్త రూపాలను సమర్థించవచ్చు. ఉదాహరణకు, కింది నిర్వచనంలోలా
రెండు ప్రమేయాలను ఒకేసారి నిర్వచించడం తరచూ ఉపయోగపడుతుంది:
\begin{align*}
h_0(\vec x, 0) & = f_0(\vec x) \\
h_1(\vec x, 0) & = f_1(\vec x) \\
h_0(\vec x, y+1) & = g_0(\vec x, y, h_0(\vec x, y), h_1(\vec x, y)) \\
h_1(\vec x, y+1) & = g_1(\vec x, y, h_0(\vec x, y), h_1(\vec x, y))
\end{align*}
ఇది \emph{సమకాల పునరావృత్తి}కి ఒక ఉదాహరణ. ప్రమేయాలను
నిర్వచించడానికి మరొక ఉపయోగకరమైన పద్ధతి, కింది నిర్వచనంలోలా
$h(\vec x,y+1)$ విలువను $h(\vec x,0)$, \dots,~$h(\vec x,y)$
\emph{విలువలన్నిటి} పరంగా ఇవ్వడం:
\begin{align*}
h(\vec x, 0) & = f(\vec x) \\
h(\vec x, y+1) & = g(\vec x, y, \tuple{h(\vec x, 0), \dots, h(\vec x, y)}).
\end{align*}
కింది నమూనా ఇదే భావాన్ని మరింత క్లుప్తంగా చూపుతుంది:
\[
h(\vec x, y) = g(\vec x, y, \tuple{h(\vec x, 0), \dots, h(\vec x, y-1)})
\]
$y$ అనేది $0$ అయినప్పుడు $g$కు చివరి ఆర్గుమెంటు కేవలం ఖాళీ క్రమమని
అర్థం చేసుకోవాలి. రెండు రూపాల్లోనూ భావం ఇదే: ``ఉత్తరగామి దశ''ను
గణించేటప్పుడు, అప్పటివరకు గణించిన విలువల మొత్తం క్రమాన్ని $h$
ఉపయోగించుకోవచ్చు. దీన్ని \emph{పూర్వవిలువల పునరావృత్తి} అంటారు.
ఒక నిర్దిష్ట ఉదాహరణగా, కింది రకపు నిర్వచనాన్ని సమర్థించడానికి దీన్ని
వాడవచ్చు:
\begin{align*}
h(\vec x, y) & = \begin{cases}
  g(\vec x, y, h(\vec x, k(\vec x, y))) & \text{$k(\vec x,y)<y$ అయితే} \\
  f(\vec x) & \text{ఇతర సందర్భాల్లో.}
\end{cases}
\end{align*}
మరో విధంగా, $k$ ఇచ్చే \emph{ఏ} పూర్వ విలువ వద్దనైనా $h$ విలువ పరంగా,
$y$ వద్ద $h$ విలువను గణించవచ్చు.


\begin{prob}
  శేష ప్రమేయం $r(x,y)$ను పూర్వవిలువల పునరావృత్తితో నిర్వచించండి.
  ($x$, $y$ సహజ సంఖ్యలు, $y>0$ అయితే, ఏదో~$z$కు
  $x=z\times y+r(x,y)$ అయ్యే $y$కన్నా చిన్న సంఖ్య $r(x,y)$.
  నిర్దిష్టత కోసం $y=0$ అయితే $r(x,0)=0$ అని తీసుకుందాం.)
\end{prob}

సాధారణ ఆదిమ పునరావృత్తిని ఉపయోగించి ఈ ప్రమేయాలను ఎలా పొందాలో
ఆలోచించాలి. ఆదిమ పునరావృత్తి యొక్క చివరి రూపం మరింత అనువైనది:
దారిలో \emph{పారామితులను} (ప్రక్క విలువలను) మార్చడానికి అనుమతిస్తుంది:
\begin{align*}
h(\vec x, 0) & = f(\vec x) \\
h(\vec x, y+1) & = g(\vec x, y, h(k(\vec x), y))
\end{align*}
దీనినీ సాధారణ ఆదిమ పునరావృత్తితో అనుకరించవచ్చు. (అలా చేయడం
సంక్లిష్టం. సూచనగా, గణనను చేతితో వెనక్కి విప్పి చూడండి.)

\end{document}
